\begin{frame}{Ausblick: Modelle weiterentwickeln}
\small
\begin{itemize}\setlength\itemsep{0.25em}
    \item \textbf{Mehr Domänendaten:} eigene Video-Frames, Innenräume, schlechte Beleuchtung und kleine Gesichter nachtrainieren.
    \item \textbf{Active Learning:} nur schwierige Fehlfälle manuell labeln: verpasste Gesichter, Profile, verdeckte Gesichter.
    \item \textbf{Temporal stabilisieren:} Detektion mit Tracking kombinieren, damit Gesichter zwischen Frames nicht flackern.
    \item \textbf{Threshold kalibrieren:} je Anwendung anderer Punkt zwischen Datenschutz-Recall und Fehlalarmrate.
    \item \textbf{Deployment optimieren:} Export nach ONNX/TensorRT, Quantisierung und Benchmark auf Zielhardware.
    \item \textbf{Privacy-QA:} automatische Stichproben nach dem Blur, besonders bei kleinen und unscharfen Gesichtern.
\end{itemize}
\end{frame}

\begin{frame}{Ausblick: mögliche Anwendungen}
\footnotesize
\begin{columns}[T,totalwidth=\textwidth]
\begin{column}{0.48\textwidth}
\textbf{Datenschutz-Pipeline}
\begin{itemize}\setlength\itemsep{0.22em}
    \item Video-Anonymisierung vor Veröffentlichung
    \item Dashcam-, Bodycam- und Eventmaterial
    \item Forschungsdatensätze mit weniger Personenbezug
\end{itemize}
\end{column}
\begin{column}{0.48\textwidth}
\textbf{Produktive Erweiterungen}
\begin{itemize}\setlength\itemsep{0.22em}
    \item Vorblur für Medienproduktion
    \item Label-Assistenz für neue Trainingsdaten
    \item Datenschutzfilter für Kameraanalyse
\end{itemize}
\end{column}
\end{columns}

\vspace{0.5em}
\scriptsize
\begin{tabular}{p{0.28\textwidth} p{0.62\textwidth}}
\toprule
Starkes Modell & Nächster Schritt \\
\midrule
Red6 ep25 & guter Spezialist, weiter mit schwierigen Fehlfällen nachtrainieren \\
Full ep30 & bessere Generalisierung, für produktive Pipeline zuerst kalibrieren \\
\bottomrule
\end{tabular}

\vspace{0.25em}
{\scriptsize Kerngedanke: Hoher Recall ist für Datenschutz wichtiger als perfekte Boxen.}
\end{frame}
